\documentclass[11pt,a4paper]{article}
\usepackage[T1]{fontenc}
\usepackage[utf8]{inputenc}
\usepackage{lmodern}
\usepackage[margin=27mm]{geometry}
\usepackage{amsmath,amssymb,amsthm,mathtools,booktabs,tabularx,microtype}
\usepackage[dvipsnames]{xcolor}
\usepackage[colorlinks=true,linkcolor=MidnightBlue,citecolor=MidnightBlue,urlcolor=MidnightBlue]{hyperref}
\hypersetup{pdftitle={Transient objective convergence of spectrally weighted gradient descent},pdfauthor={Anonymous}}
\newtheorem{theorem}{Theorem}
\newtheorem{proposition}[theorem]{Proposition}
\newtheorem{lemma}[theorem]{Lemma}
\newtheorem{corollary}[theorem]{Corollary}
\theoremstyle{remark}
\newtheorem{remark}[theorem]{Remark}
\newcommand{\R}{\mathbb R}
\newcommand{\diag}{\operatorname{diag}}
\newcommand{\E}{\mathbb E}
\newcommand{\W}{\mathcal W}
\newcommand{\Ck}{C_\kappa}
\title{Transient objective convergence of spectrally weighted gradient descent}
\author{Anonymous}
\date{6 November 2026\\\small Version 0.1.0-candidate\\\href{https://doi.org/10.5281/zenodo.24001734}{doi:10.5281/zenodo.24001734}}

\begin{document}
\maketitle
\begin{center}\small AI-assisted research preprint; not externally peer reviewed.\end{center}

\begin{abstract}
Spectrally weighted gradient methods can share the long-time asymptotic rate of steepest descent while having a different worst-case objective decrease over finite horizons. We study this distinction for positive definite quadratic objectives normalised by $\lambda_{\min}(A)=1$ and $\lambda_{\max}(A)=\kappa$. This normalisation is part of the statement for a general nonhomogeneous weight. For every positive, nondecreasing, unbounded spectral weight, every fixed number $K$ of iterations, and every fixed relaxation parameter $0<\omega\le1$, the worst-case objective ratio has large-condition-number expansion $1-2\omega K/\kappa+o(\kappa^{-1})$. An elementary upper bound gives one direction; two- and three-dimensional constructions give the matching lower bound. For power weights $\Psi(\lambda)=\lambda^{2\ell}$ with $\ell>0$, the unrelaxed result remains sharp for every horizon $K=o(\log\kappa)$. A quantitative three-mode estimate also produces simultaneous failures of the per-iteration Kantorovich bound over logarithmically many iterations. Exact rational examples refute the two-step bound already for the standard minimal-gradient method. Two active eigenvalues cannot produce this failure: their two-step map is scalar and always satisfies the Kantorovich bound. These results answer a strong interpretation of a recent two-step question; we also prove the weaker single-factor interpretation and distinguish both from the established fixed-matrix asymptotic theory.
\end{abstract}

\section{Introduction}

Consider the quadratic minimisation problem
\begin{equation}\label{eq:objective}
 \min_{x\in\R^n} f(x),\qquad
 f(x)=\tfrac12x^TAx-b^Tx,\qquad A=A^T\succ0.
\end{equation}
Its minimiser is $x_*=A^{-1}b$. We study gradient-direction methods
\begin{equation}\label{eq:iteration}
 x_{j+1}=x_j-\omega\alpha_jg_j,\qquad
 g_j=Ax_j-b,\qquad
 \alpha_j=\frac{g_j^T\Psi(A)g_j}{g_j^TA\Psi(A)g_j}.
\end{equation}
For $\Psi\equiv1$ and $\omega=1$, this is steepest descent with exact line search. For $\Psi(\lambda)=\lambda$, it is the minimal-gradient method, which minimises the Euclidean norm of the next gradient along the current gradient direction. The family $\Psi(\lambda)=\lambda^{2\ell}$ includes both cases and corresponds to $\ell$-minimal gradient descent in \cite{MMT2025}. The notation of Pronzato, Wynn and Zhigljavsky \cite{PWZ2006} corresponds to $P(\lambda)=\Psi(\lambda)/\lambda$; their $P$-gradient family imposes additional regularity conditions on the weight.

Let
\[
 E_j=f(x_j)-f(x_*)=\tfrac12(x_j-x_*)^TA(x_j-x_*),
 \qquad \kappa=\lambda_{\max}(A)/\lambda_{\min}(A).
\]
The ordinary objective gap $E_j$ is the quantity studied throughout. Results for the norm minimised by the step rule need not transfer to this objective without a condition-number factor. The mechanism is simple: a high-curvature component can carry little objective energy yet dominate the weighted choice of step. The algorithm then takes short steps while most objective error remains in a low-curvature component. Two high modes keep this imbalance alive even without relaxation.

For steepest descent, the classical Kantorovich estimate is
\begin{equation}\label{eq:kantorovich}
 E_{j+1}\le\Ck E_j,\qquad
 \Ck=\left(\frac{\kappa-1}{\kappa+1}\right)^2.
\end{equation}
Pronzato et al.\ establish common long-time asymptotic behaviour for a broad family of $P$-gradient methods \cite{PWZ2006}. Huang et al.\ give explicit limiting ratios for odd and even subsequences \cite{HDLZ2022}. In particular, Corollary 11, equations (60)--(62), and Remark 12, equation (63), of their author manuscript arXiv:1908.07111v1 imply that the limiting two-step objective ratio is at most $\Ck^2$, under their hypotheses. Those statements first fix the matrix and starting point and then send the iteration count to infinity. For symmetric positive-definite systems, standard minimal gradient is also known as minimal residual iteration or GMRES(1). He \cite{He2025} studies its residual-norm and root-convergence factors; these differ from the ordinary-objective finite-horizon quantity considered here.

MacDonald, Murray and Tappenden \cite{MMT2025} develop nonasymptotic results for relaxed power-weighted methods. Their published Section 6, paragraph headed ``Future work'', asks whether a two-step objective decrease can attain the Kantorovich rate. The question is expressed in prose, rather than as a quantified inequality. Two interpretations must therefore be separated:
\begin{align}
 E_{j+2}&\le\Ck^2E_j, &&\text{the per-iteration rate over two steps},\label{eq:strong}\\
 E_{j+2}&\le\Ck E_j, &&\text{one Kantorovich factor over two steps}.\label{eq:weak}
\end{align}
We disprove \eqref{eq:strong}, including for the standard minimal-gradient method, and prove \eqref{eq:weak} in the unrelaxed case. We do not attribute the fully quantified statement \eqref{eq:strong} verbatim to the source.

The later flexible framework of MacDonald, Murray and Tappenden \cite[Definition 8 and Section 2.1, Lemma 9]{MMT2025b} permits multiple search directions and symmetric preconditioning, with comparison in a selected weighted gradient norm. Those results do not by themselves give the sharp ordinary-objective transient lower construction studied here. Our counterexamples concern the single-gradient-direction update \eqref{eq:iteration}, not every algorithm in that broader framework.

The main result identifies the sharp leading coefficient in the finite-horizon worst-case objective decrease. In the unrelaxed case $\omega=1$, this coefficient is $2K$ for every positive, nondecreasing, unbounded spectral weight, whereas steepest descent has coefficient $4K$. Power-weighted methods exhibit this distinction over horizons growing more slowly than $\log\kappa$, and explicit instances violate the Kantorovich bound at every horizon in an interval of length proportional to $\log\kappa$.

The methods and their conjugacy to steepest descent are established objects. Our proposed contribution is the objective-specific finite-horizon lower construction, its quantitative control, and the resulting order-of-limits distinction. No claim is made that the classical two-point dynamics or the underlying algorithm family is new.

\subsection*{Reading the result: metrics and limits}
\begin{center}\small
\begin{tabularx}{\textwidth}{@{}p{.23\textwidth}XX@{}}
\toprule
Quantity or regime & What is held fixed & Scope\\
\midrule
Ordinary objective $E_j$ & The displayed quadratic & Measured error, not the weighted gradient norm choosing the step.\\
Fixed horizon (Theorem \ref{thm:fixed}) & $\Psi,\omega,K$; $\kappa\to\infty$ & Coefficient $2\omega K$; instances may vary with $\kappa$.\\
Sublogarithmic (Theorem \ref{thm:moving}) & $\ell>0$, $\omega=1$; $K=o(\log\kappa)$ & Coefficient two per step; witness may depend on the chosen horizon.\\
Simultaneous (Corollary \ref{cor:log}) & $\ell>0$, admissible $\theta$ & One instance at each $\kappa$ works for every indicated horizon.\\
Fixed-matrix long time & One matrix and starting point; $j\to\infty$ & Inherited asymptotic theory, a different limit.\\
\bottomrule
\end{tabularx}
\end{center}
\paragraph{Proof roadmap.}
Proposition \ref{prop:upper} gives an objective upper bound. Theorem \ref{thm:fixed} matches its leading deficit using a small amount of high-mode energy. Proposition \ref{prop:two} explains why two active eigenvalues cannot give the unrelaxed two-step failure. Lemma \ref{lem:uniform} controls the perturbation for finite parameters, giving an explicit sufficient failure test and then the growing-horizon results. Readers seeking a concrete check can first consult the exact rational examples.

\section{Normalisation and an objective bound}

Translation by $x_*$ reduces \eqref{eq:objective} to $b=0$. We impose
\begin{equation}\label{eq:normalise}
 \lambda_{\min}(A)=1,\qquad \lambda_{\max}(A)=\kappa>1.
\end{equation}
For power weights this loses no generality after rescaling $A$. For a general nonhomogeneous $\Psi$, the normalisation is part of the statement. A rescaling may change the function used as a spectral weight.

Unless stated otherwise, $\Psi:[1,\infty)\to(0,\infty)$ is nondecreasing, and $0<\omega\le1$. Finite-dimensional spectral calculus defines $\Psi(A)$ without continuity assumptions. Multiplying $\Psi$ by a positive scalar does not change the algorithm. When an iterate reaches the minimiser, every later iterate is defined to equal it.

\begin{proposition}\label{prop:upper}
Under these assumptions, every iterate satisfies
\begin{equation}\label{eq:upper}
 E_{j+1}\le(1-\omega/\kappa)^2 E_j,
 \qquad E_K\le(1-\omega/\kappa)^{2K}E_0.
\end{equation}
\end{proposition}
\begin{proof}
Write $\alpha_{\rm SD}=g^Tg/(g^TAg)$. The reciprocal weighted-mean formula and monotonicity of $\Psi$ imply
\begin{equation}\label{eq:stepinterval}
 \kappa^{-1}\le\alpha\le\alpha_{\rm SD}.
\end{equation}
For completeness, assign probability weights proportional to the squared components of $g$ in an eigenbasis, and let $\Lambda$ be the resulting spectral random variable. If $\Lambda'$ is an independent copy, then
\[
 \E[\Lambda\Psi(\Lambda)]-\E[\Lambda]\E[\Psi(\Lambda)]
 =\tfrac12\E[(\Lambda-\Lambda')(\Psi(\Lambda)-\Psi(\Lambda'))]\ge0.
\]
This proves the second inequality in \eqref{eq:stepinterval}; the first follows from $\Lambda\le\kappa$.

The quadratic $t\mapsto f(x-tg)$ decreases for $0\le t\le\alpha_{\rm SD}$. Both $\omega/\kappa$ and $\omega\alpha$ lie in this interval, in the stated order. Hence
\[
 f(x-\omega\alpha g)\le f(x-\omega g/\kappa).
\]
In an eigenbasis, the last update multiplies the coordinate with eigenvalue $\lambda$ by $1-\omega\lambda/\kappa\in[0,1-\omega/\kappa]$. This proves the one-step estimate. Iteration proves the second.
\end{proof}

\begin{corollary}\label{cor:weak}
For $\omega=1$, every nondecreasing positive weight satisfies \eqref{eq:weak}.
\end{corollary}
\begin{proof}
Proposition \ref{prop:upper} gives $E_{j+2}/E_j\le(1-1/\kappa)^4$. Moreover
\[
 (1-1/\kappa)^4\le\left(\frac{\kappa-1}{\kappa+1}\right)^2
\]
because $(\kappa-1)(\kappa+1)\le\kappa^2$.
\end{proof}

The comparison argument in Proposition \ref{prop:upper} is elementary. It improves the Rate I bound of \cite[Theorem 10]{MMT2025} in the range $0<\omega\le1$, but its independent historical priority is not needed for the lower-bound results below.

\section{Sharp fixed-horizon asymptotics}

For an integer $K\ge1$, define
\begin{equation}\label{eq:worstcase}
 \W_{\Psi,\omega,K}(\kappa)
 =\sup\frac{E_K}{E_0},
\end{equation}
where the supremum ranges over finite dimensions, matrices satisfying \eqref{eq:normalise}, and starting points with $E_0>0$. Let $\W^{(n)}$ denote the same supremum restricted to dimension $n$.

\begin{theorem}\label{thm:fixed}
Suppose that $\Psi$ is positive, nondecreasing and unbounded. For each fixed $0<\omega\le1$ and positive integer $K$,
\begin{equation}\label{eq:fixed}
 \lim_{\kappa\to\infty}\kappa[1-\W_{\Psi,\omega,K}(\kappa)]
 =2\omega K.
\end{equation}
The same limit holds with $\W$ replaced by $\W^{(3)}$. If $0<\omega<1$, it also holds with $\W^{(2)}$.
\end{theorem}

\begin{proof}
Proposition \ref{prop:upper} yields
\[
 \liminf_{\kappa\to\infty}\kappa(1-\W_{\Psi,\omega,K})\ge2\omega K.
\]
It remains to construct matching examples. We describe an initial point by its objective-energy weights $w_i=\lambda_i x_{0,i}^2$, so that $2E_0=\sum_iw_i$. Write
\[
 P_j(z)=\prod_{i=0}^{j-1}(1-\omega\alpha_i z).
\]
Then $2E_j=\sum_iw_iP_j(\lambda_i)^2$.

First take $0<\omega<1$, let $A_\kappa=\diag(1,\kappa)$, and set
\[
 T_\kappa=\sqrt{\Psi(\kappa)},\qquad
 (w_1,w_2)=\left(1,\frac{T_\kappa}{\kappa\Psi(\kappa)}\right).
\]
The high-mode weight is $o(\kappa^{-1})$. Put $H_j=P_j(1)^2$, $V_j=P_j(\kappa)^2$, $a_j=\kappa\alpha_j$, and $\varepsilon_\kappa=\Psi(1)/T_\kappa\to0$. The exact recurrence is
\[
 a_j=\frac{\varepsilon_\kappa H_j+V_j}{\varepsilon_\kappa H_j/\kappa+V_j},
 \quad H_{j+1}=H_j(1-\omega a_j/\kappa)^2,
 \quad V_{j+1}=V_j(1-\omega a_j)^2.
\]
For each fixed $j$, induction gives $a_j\to1$, $H_j\to1$, and $V_j\to(1-\omega)^{2j}>0$. Therefore
\[
 H_K=1-\frac{2\omega K}{\kappa}+o(\kappa^{-1}),\qquad
 \frac{E_K}{E_0}=H_K+o(\kappa^{-1}),
\]
which proves the matching upper limit for the scaled deficit.

Now take $\omega=1$. Fix $c\in(0,1)$ and require $\kappa\ge1/c$. Set
\begin{equation}\label{eq:generalconstruction}
 A_\kappa=\diag(1,c\kappa,\kappa),\quad
 T_\kappa=\sqrt{\Psi(c\kappa)},\quad
 (w_1,w_2,w_3)=\left(1,\frac{T_\kappa}{c\kappa\Psi(c\kappa)},
 \frac{T_\kappa}{\kappa\Psi(\kappa)}\right).
\end{equation}
The total high-mode energy is at most $(1+c^{-1})/[\kappa\sqrt{\Psi(c\kappa)}]=o(\kappa^{-1})$. Its two contributions to the step numerator initially equal $T_\kappa$. Put $U_j=P_j(c\kappa)^2$, $V_j=P_j(\kappa)^2$, $H_j=P_j(1)^2$, $a_j=\kappa\alpha_j$, and $\varepsilon_\kappa=\Psi(1)/T_\kappa$. Then
\begin{align}
 a_j&=\frac{U_j+V_j+\varepsilon_\kappa H_j}
 {cU_j+V_j+\varepsilon_\kappa H_j/\kappa},\label{eq:recurrence}\\
 H_{j+1}&=H_j(1-a_j/\kappa)^2,\quad
 U_{j+1}=U_j(1-ca_j)^2,\quad
 V_{j+1}=V_j(1-a_j)^2.\notag
\end{align}
Let $d=(1-c)/(1+c)$. Since $\varepsilon_\kappa\to0$, induction for each fixed $j$ gives
\begin{equation}\label{eq:limiting}
 a_j\longrightarrow\frac2{1+c},\qquad H_j\longrightarrow1,
 \qquad U_j,V_j\longrightarrow d^{2j}>0.
\end{equation}
The induction is valid because every limiting denominator is positive. At equal high-mode weights the limiting step is $2/(1+c)$, and both high-mode squared multipliers equal $d^2$.

Expanding the finite product for $H_K$ and using the negligible high-mode energy gives
\begin{equation}\label{eq:constructionasymp}
 \frac{E_K}{E_0}=1-\frac{4K}{(1+c)\kappa}+o(\kappa^{-1}).
\end{equation}
Thus $\limsup\kappa(1-\W^{(3)})\le4K/(1+c)$ for each $c<1$. Letting $c\uparrow1$ proves \eqref{eq:fixed}. The two-dimensional construction for $\omega<1$ embeds in dimension three by adjoining a zero initial coordinate at an existing eigenvalue.
\end{proof}

\begin{corollary}\label{cor:block}
Fix an unbounded positive nondecreasing $\Psi$, $0<\omega\le1$, and an integer $K\ge1$. The inequality
\begin{equation}\label{eq:relaxedkant}
 E_K\le\left(1-\omega(2-\omega)\frac{4\kappa}{(1+\kappa)^2}\right)^K E_0
\end{equation}
does not hold uniformly over positive definite quadratic problems.
\end{corollary}
\begin{proof}
The factor on the right has expansion $1-4\omega(2-\omega)K/\kappa+O(\kappa^{-2})$. Its leading deficit exceeds $2\omega K$, because $0<\omega\le1$. Theorem \ref{thm:fixed} gives examples with a strictly larger objective ratio for all sufficiently large condition numbers.
\end{proof}

For $\Psi(\lambda)=\lambda^{2\ell}$, every fixed $\ell>0$ satisfies Theorem \ref{thm:fixed}. At $\ell=0$, the unrelaxed exact worst-case factor is $\Ck^K$, whose leading deficit is $4K/\kappa$. Hence the limit $\ell\downarrow0$ is not uniform in $\kappa$. This statement concerns the worst objective ratio; it does not assert a discontinuity of individual iterates at fixed $A$.

\section{Why three spectral modes are necessary}

The two-point dynamics explain both the limiting construction and why two-dimensional tests can miss the obstruction. The following identity is a direct form of the classical two-point recurrence underlying \cite{PWZ2006,HDLZ2022}; see also \cite[Theorem 9 and Lemma 10]{He2025} for minimal residual iteration. We include its proof to make the dimensional boundary explicit.

\begin{proposition}\label{prop:two}
Suppose $\omega=1$, $\Psi$ is positive, and the initial error has components in at most two distinct eigenspaces of $A$. Then
\begin{equation}\label{eq:two}
 E_{j+2}\le\Ck^2 E_j
\end{equation}
for every $j$, and every even block satisfies $E_{j+2m}\le\Ck^{2m}E_j$.
\end{proposition}
\begin{proof}
If only one active eigenvalue remains, one step reaches the minimiser. Otherwise let $a$ be the lower active eigenvalue and let the upper active eigenvalue be $a\rho$, where $1<\rho\le\kappa$. Rescale the matrix to $A/a$ and define the rescaled weight $\widetilde\Psi(t)=\Psi(at)$. The gradient is divided by $a$ and the stepsize is multiplied by $a$, leaving the iterate update unchanged. At a given iterate define
\[
 R=\frac{\widetilde\Psi(\rho)g_{\rho}^2}{\widetilde\Psi(1)g_1^2}>0,
\]
where $g_1^2$ and $g_\rho^2$ denote squared Euclidean norms of the eigenspace projections of the rescaled gradient. The rescaled stepsize is $\widetilde\alpha=(1+R)/(1+\rho R)$. Its multipliers on the low and high components are respectively
\[
 \frac{R(\rho-1)}{1+\rho R},\qquad
 -\frac{\rho-1}{1+\rho R}.
\]
Consequently the next weighted ratio is $R^{-1}$. After two steps, both error components are multiplied by the same positive scalar
\[
 \gamma(R)=\frac{(\rho-1)^2R}{(1+\rho R)(\rho+R)}
 =\frac{(\rho-1)^2}{\rho(R+R^{-1})+\rho^2+1}
 \le\left(\frac{\rho-1}{\rho+1}\right)^2.
\]
The energy ratio over these two steps is $\gamma(R)^2$. Since $\rho\le\kappa$, this proves \eqref{eq:two}; iteration proves the even-block estimate.
\end{proof}

Three active eigenvalues are therefore the smallest possible support of an unrelaxed two-step counterexample. This is a statement about the active spectrum, so repeated eigenvalues or additional zero coordinates do not alter it.

\section{Quantitative control and growing horizons}

We now take $\Psi(\lambda)=\lambda^{2\ell}$ with fixed $\ell>0$, and $\omega=1$. Write $\W_{\ell,K}$ for the corresponding worst-case ratio. The finite-horizon construction can be controlled while $K$ grows with $\kappa$.

Fix
\begin{equation}\label{eq:parameters}
 \kappa\ge4,\quad \tfrac12\le c<1,\quad
 \delta=1-c,\quad d=\frac{1-c}{1+c},\quad
 0<\nu<2\ell,\quad\mu=2\ell-\nu.
\end{equation}
Use the three-mode matrix and starting point
\begin{equation}\label{eq:powerconstruction}
 A_\kappa=\diag(1,c\kappa,\kappa),\qquad
 x_0=\left(1,c^{-(\ell+1)}\kappa^{-(\nu+2)/2},
 \kappa^{-(\nu+2)/2}\right).
\end{equation}
Its energy weights are $(1,c^{-(2\ell+1)}\kappa^{-1-\nu},\kappa^{-1-\nu})$. Set
\[
 \beta_j=\kappa^{-\mu}d^{-2j},\qquad C_0=1+c^{-(2\ell+1)}.
\]

\begin{lemma}\label{lem:uniform}
If $K\ge1$ and $\beta_K\le\delta/64$, then the construction \eqref{eq:powerconstruction} satisfies
\begin{equation}\label{eq:quantitative}
 \left|\kappa\left(1-\frac{E_K}{E_0}\right)-\frac{4K}{1+c}\right|
 \le\beta_K+C_0\kappa^{-\nu}+\frac{9K^2}{2\kappa}.
\end{equation}
\end{lemma}

\begin{proof}
The recurrence \eqref{eq:recurrence} holds with $\varepsilon_\kappa=\kappa^{-\mu}$. Whenever the high-mode multipliers are nonzero, define
\[
 r_j=\log(U_j/V_j),\qquad s_j=\log(V_j/d^{2j}),\qquad
 \epsilon_j=\kappa^{-\mu}H_j/V_j=\beta_jH_je^{-s_j}.
\]
The variables $r_j$ and $s_j$ measure, respectively, the high-mode imbalance and departure from their common limiting decay. Since $\alpha_j$ is the reciprocal of a weighted mean of eigenvalues at least one, $0<\alpha_j\le1$ and hence $0\le H_j\le1$. We prove simultaneously that
\begin{equation}\label{eq:bootstrap}
 |r_j|,|s_j|\le4\beta_j/\delta,\qquad\epsilon_j\le2\beta_j
 \quad(0\le j\le K).
\end{equation}
At $j=0$ the bounds hold. Under the stated hypothesis, the claimed logarithmic bounds do not exceed $1/16$.

Put $A_c=(c-1/\kappa)/\delta$ and $B_c=(1-1/\kappa)/\delta$. Both are nonnegative and at most $1/\delta$. In terms of $r_j,\epsilon_j$,
\begin{equation}\label{eq:eta}
 a_j=\frac{e^{r_j}+1+\epsilon_j}{ce^{r_j}+1+\epsilon_j/\kappa}.
\end{equation}
Direct algebra gives
\[
 1-ca_j=\frac{\delta-(c-1/\kappa)\epsilon_j}
 {ce^{r_j}+1+\epsilon_j/\kappa},\quad
 1-a_j=-\frac{\delta e^{r_j}+(1-1/\kappa)\epsilon_j}
 {ce^{r_j}+1+\epsilon_j/\kappa}.
\]
The first numerator remains positive because $A_c\epsilon_j\le1/32$; the second is strictly positive before its minus sign. Thus neither high mode vanishes at the next step. Their logarithmic recurrences are
\begin{align}
 r_{j+1}&=-r_j+2\log(1-A_c\epsilon_j)
 -2\log(1+B_c\epsilon_je^{-r_j}),\label{eq:rlog}\\
 s_{j+1}-s_j&=2r_j-2\log\frac{ce^{r_j}+1}{1+c}
 +2\log(1+B_c\epsilon_je^{-r_j})\notag\\
 &\hspace{13mm}-2\log\left(1+\frac{\epsilon_j}{\kappa(ce^{r_j}+1)}\right).
 \label{eq:slog}
\end{align}
For $0\le u\le1/32$, $-\log(1-u)\le4u/3$; for $v\ge0$, $\log(1+v)\le v$. Using $e^{1/16}\le16/15$ also makes these numerical constant comparisons rational. Hence
\begin{equation}\label{eq:rerror}
 |r_{j+1}+r_j|<5\epsilon_j/\delta\le10\beta_j/\delta.
\end{equation}
The derivative of $2r-2\log[(ce^r+1)/(1+c)]$ is $2/(ce^r+1)\le2$. The remaining terms in \eqref{eq:slog} have total absolute value at most
\[
 2e^{1/16}\epsilon_j/\delta+2\epsilon_j/\kappa
 \le(2e^{1/16}+1/4)\epsilon_j/\delta<3\epsilon_j/\delta.
\]
Therefore
\begin{equation}\label{eq:serror}
 |s_{j+1}-s_j|\le2|r_j|+3\epsilon_j/\delta
 \le14\beta_j/\delta.
\end{equation}
As $d\le1/3$,
\begin{equation}\label{eq:betasum}
 \sum_{i=0}^{j-1}\beta_i
 =\beta_j\frac{d^2(1-d^{2j})}{1-d^2}\le\beta_j/8.
\end{equation}
Summing \eqref{eq:rerror} in absolute value and \eqref{eq:serror} gives respectively $|r_j|\le5\beta_j/(4\delta)$ and $|s_j|\le7\beta_j/(4\delta)$, strictly within \eqref{eq:bootstrap}. Finally $\epsilon_j\le e^{1/16}\beta_j<2\beta_j$. This closes the induction, including nonvanishing of the high modes.

Compare first $\epsilon_j$ to zero at fixed $r_j$, along the nonnegative $\epsilon$ interval, and then compare $r_j$ to zero at $\epsilon=0$. The bootstrap keeps $|r|\le1/16$ on this path. At $\epsilon=0$, the absolute derivative in $r$ of the map in \eqref{eq:eta} equals $\delta e^r/(ce^r+1)^2\le\delta/(4c)\le\delta/2$. Its derivative in $\epsilon$ is nonnegative and at most one, because
\[
 0\le\frac{ce^r+1-(e^r+1)/\kappa}
 {(ce^r+1+\epsilon/\kappa)^2}\le1.
\]
It follows that
\[
 |a_j-2/(1+c)|\le4\beta_j,\qquad
 2\sum_{j<K}|a_j-2/(1+c)|\le\beta_K.
\]
Also $a_j\le4/3+1/32<3/2$. Put $z=\sum_{j<K}a_j/\kappa$. The elementary product and exponential bounds imply
\[
 1-2z\le H_K=\prod_{j<K}(1-a_j/\kappa)^2
 \le e^{-2z}\le1-2z+2z^2.
\]
Consequently
\[
 \left|\kappa(1-H_K)-\frac{4K}{1+c}\right|
 \le\beta_K+\frac{9K^2}{2\kappa}.
\]
For $K\ge1$, the bootstrap gives $V_K\le e^{1/16}d^{2K}<1$ and $U_K\le e^{1/8}d^{2K}<1$. The exact objective ratio is
\[
 \frac{E_K}{E_0}
 =\frac{H_K+\kappa^{-1-\nu}[c^{-(2\ell+1)}U_K+V_K]}
 {1+C_0\kappa^{-1-\nu}},
\]
so $\kappa|E_K/E_0-H_K|\le C_0\kappa^{-\nu}$. Combining the last two estimates proves \eqref{eq:quantitative}.
\end{proof}

\subsection*{A finite-parameter applicability and failure test}
Lemma \ref{lem:uniform} can be used without sending $\kappa$ to infinity. Retain all parameters and the initial point in \eqref{eq:parameters}--\eqref{eq:powerconstruction}, and put
\[
 B_K=\kappa^{-\mu}d^{-2K}+C_0\kappa^{-\nu}+\frac{9K^2}{2\kappa}.
\]
Its hypothesis on the integer horizon is exactly
\[
 1\le K\le\left\lfloor
 \frac{\mu\log\kappa+\log(\delta/64)}{2\log(1/d)}
 \right\rfloor.
\]
If the upper endpoint is less than one, this test certifies no positive horizon. For any admitted $K$, the strict inequality
\[
 \frac{4K}{1+c}+B_K<\kappa\left[1-\Ck^K\right]
\]
is sufficient to certify $E_K/E_0>\Ck^K$ for that specific initial point.

\paragraph{Derivation.}
Taking logarithms of $\kappa^{-\mu}d^{-2K}\le\delta/64$ and using $\log(1/d)>0$ gives the horizon condition. Lemma \ref{lem:uniform} gives $E_K/E_0\ge1-[4K/(1+c)+B_K]/\kappa$, which proves the strict failure conclusion. This is an elementary consequence of the lemma, not an optimal horizon formula or a necessary condition for failure. Numerical evaluation of nonrational powers or logarithms requires certified enclosures if it is to count as a certificate; decimal evaluation alone is not an exact test. The exact rational counterexamples below establish failure directly and need not satisfy this sufficient test.

\begin{theorem}\label{thm:moving}
Fix $\ell>0$. For every integer-valued horizon $K=K(\kappa)\ge1$ with $K=o(\log\kappa)$,
\begin{equation}\label{eq:moving}
 \lim_{\kappa\to\infty}\frac{\kappa}{K(\kappa)}
 [1-\W_{\ell,K(\kappa)}(\kappa)]=2.
\end{equation}
The same statement holds when the supremum is restricted to dimension three.
\end{theorem}
\begin{proof}
Proposition \ref{prop:upper} gives the lower limit two for the scaled deficit, since $K=o(\kappa)$. For a matching example, take $\nu=\ell$ in Lemma \ref{lem:uniform}, write $h=\log\kappa$, and choose
\[
 d=\exp[-\sqrt{h/K}],\qquad c=(1-d)/(1+d).
\]
Then $d\to0$, $c\to1$, and eventually $c\ge1/2$. Since $\delta=2d/(1+d)\ge d$,
\[
 \frac{\beta_K}{\delta}
 \le\exp[-\ell h+2\sqrt{Kh}+\sqrt{h/K}]\longrightarrow0.
\]
The exponent tends to minus infinity because $K/h\to0$ and $K\ge1$. Thus Lemma \ref{lem:uniform} applies eventually. Its error divided by $K$ tends to zero, whereas $4/(1+c)=2(1+d)\to2$. One admissible three-dimensional instance lower-bounds the worst-case ratio and gives the required upper limit for its deficit.
\end{proof}

\begin{corollary}\label{cor:log}
Fix $\ell>0$ and $0<\theta<2\ell/\log9$. There exists a three-dimensional family such that, for all sufficiently large $\kappa$,
\begin{equation}\label{eq:logfailure}
 \frac{E_K}{E_0}>\Ck^K
 \qquad\text{for every integer }1\le K\le\theta\log\kappa.
\end{equation}
The same instance at each $\kappa$ works simultaneously for every indicated horizon.
\end{corollary}
\begin{proof}
Take $c=1/2$, so $d=1/3$, and choose $0<\nu<2\ell-\theta\log9$. Then $\mu=2\ell-\nu>\theta\log9$, and $\beta_K\le\kappa^{-\mu+\theta\log9}\to0$ uniformly on the stated range. Lemma \ref{lem:uniform} yields
\[
 \kappa(1-E_K/E_0)=\frac{8K}{3}+o(1)
\]
uniformly on that range. In contrast,
\[
 \kappa(1-\Ck^K)=4K+O(K^2/\kappa+K/\kappa),
\]
with a uniform error tending to zero. The positive leading difference $4K/3$ proves \eqref{eq:logfailure}.
\end{proof}

\begin{remark}
Neither Theorem \ref{thm:moving} nor Corollary \ref{cor:log} fixes one matrix and then sends $K$ to infinity. The matrix and starting point vary with $\kappa$; in Theorem \ref{thm:moving}, their high-mode separation also depends on the chosen horizon. The results therefore coexist with the fixed-matrix long-time asymptotic theory in \cite{PWZ2006,HDLZ2022}. No sharp worst-case formula is asserted here for arbitrary horizons comparable to or larger than $\log\kappa$.
\end{remark}

\section{Exact rational counterexamples}

The limiting arguments do not depend on numerical optimisation. Small exact instances also establish the two-step failure directly.

\subsection{Standard minimal gradient}

Take $\ell=1/2$, $\omega=1$, and
\begin{equation}\label{eq:mgexample}
 A=\diag(1,60,100),\qquad x_0=(1,1/100,1/250)^T.
\end{equation}
Direct rational arithmetic gives
\[
 \alpha_0=\frac{193}{14485},\qquad
 \alpha_1=\frac{7013813}{446629608},
\]
and
\begin{equation}\label{eq:mgratio}
 \frac{E_2}{E_0}
 =\frac{830774224427957654527687}{887449697831729589460572}
 \approx0.9361366919812527.
\end{equation}
The Kantorovich two-step value is $(99/101)^4\approx0.9231138845986189$. Their exact difference is positive:
\[
 \frac{E_2}{E_0}-\left(\frac{99}{101}\right)^4
 =\frac{1202635053150759637694223270715}
 {92348371423698611602832496009372}>0.
\]
This example uses the classical minimal-gradient method itself.

\subsection{A smaller-condition-number example}

For $\ell=1$ and $\omega=1$, take
\[
 A=\diag(1,6,10),\qquad x_0=(1,1/10,1/25)^T.
\]
Then
\[
 \alpha_0=\frac{749}{5969},\qquad
 \alpha_1=\frac{285369}{1748690},\qquad
 \frac{E_2}{E_0}=\frac{145981491016797622500}{293076563677580674949}
 \approx0.4981001864666148.
\]
Here $(9/11)^4\approx0.4481251280650229$, and
\[
 \frac{E_2}{E_0}-\left(\frac9{11}\right)^4
 =\frac{214439675688327182682111}{4290933968803458661928309}>0.
\]

\subsection{A rational family for arbitrary fixed horizons}

In the unrelaxed power construction take $\ell=1/2$, $c=1/4$, $\nu=1/2$, and $\kappa=N^4$ for an integer $N\ge2$. This gives
\begin{equation}\label{eq:rationalfamily}
 A_N=\diag(1,N^4/4,N^4),\qquad
 x_0=(1,8/N^5,1/N^5)^T.
\end{equation}
Every entry is rational, and $A_N$ has integer entries when $N$ is even. The fixed-step induction in Theorem \ref{thm:fixed}, with the power-specific energy weights, gives
\begin{equation}\label{eq:rationalasymp}
 \lim_{N\to\infty}N^4\left(1-\frac{E_K}{E_0}\right)=\frac{16K}{5}
 \qquad(K\text{ fixed}).
\end{equation}
For example, $N=4$ gives $\kappa=256$ and $E_2/E_0\approx0.9711072425654474>(255/257)^4\approx0.9692330804194722$. The rational family supplies explicit finite instances for every fixed block length, although the existence theorem does not assert a small sufficient $N$ uniformly in $K$.

\section{Interpretation and limits}

The construction separates the spectrum that controls the step from the spectrum that carries most of the objective. A small amount of objective energy at large eigenvalues dominates the weighted step numerator. Two nearby high modes then maintain a short step for many iterations. Meanwhile the low mode carries almost all the ordinary objective and decreases at the rate associated with a step near $1/\kappa$.

For power weights, the change of variables $y=A^\ell(x-x_*)$ transforms the iteration into steepest descent for $\tfrac12y^TAy$. The measured objective becomes $\tfrac12y^TA^{1-2\ell}y$. This conjugacy explains why familiar weighted-norm convergence and the slower finite-time objective decrease can occur together. It does not alone establish the sharp finite-horizon coefficient or the quantitative interval of failure; those require the energy-specific construction.

The two-step issue has a precise boundary. One Kantorovich factor over two unrelaxed steps is always valid under the monotone-weight assumption. Two factors over the same steps fail with three active spectral modes, including for standard minimal gradient. No fixed block length repairs the per-iteration rate uniformly over condition number. For power weights, the sharp coefficient remains two throughout every sublogarithmic horizon, and one instance can violate the claimed per-iteration bound throughout a logarithmic interval.

The results do not give an exact finite-$\kappa$ formula for $\W_{\Psi,\omega,K}$, a sharp transition at horizons of order $\log\kappa$, or a complexity lower bound for arbitrary algorithms. They cover $0<\omega\le1$ for the fixed-horizon theorem and $\omega=1$ for the moving-horizon result. The case $1<\omega<2$ requires a separate comparison argument. The one-step upper estimate and the two-point identity are supporting ingredients whose novelty is not asserted. In particular, the factor-two difference in the leading transient decrement is not a claim that twice as many iterations are required for a fixed accuracy: the studied relative decrease tends to zero, and a fixed-fraction reduction lies outside these horizon regimes.

\section*{Verification and research status}

This is an AI-assisted research preprint, not external independent peer review. The mathematical claims rest on the proofs above. The displayed rational counterexamples can be checked with exact integer fractions; finite numerical grids do not prove the universal estimates. No proof-assistant formalisation or independent external reproduction is claimed. The accompanying verification record distinguishes current executable checks from historical provenance. The bounded prior-art comparison includes the originating question, the flexible-framework follow-up, the cited P-gradient and two-point asymptotic results, and equivalent names for minimal-gradient iteration. It did not identify the same objective-specific transient theorem, but does not establish exhaustive historical priority.

\section*{Declarations}
\paragraph{Data and code.}
The mathematical data are the explicit matrices, initial points and formulas in this paper. Accompanying checker sources and verification records document executable assurance separately.
\paragraph{Authorship and AI assistance.}
Anonymous scholarly attribution. AI systems assisted drafting, mathematical analysis, revision and internal review; this does not constitute external peer review.
\paragraph{Ethics.}
No human participants, personal data or animal research are involved.
\paragraph{Funding and competing interests.}
No funding or competing-interest declaration was supplied with the research bundle.

\begin{thebibliography}{9}

\bibitem{MMT2025}
L. MacDonald, R. Murray, and R. Tappenden,
On a family of relaxed gradient descent methods for strictly convex quadratic minimization,
\emph{Computational Optimization and Applications} \textbf{91} (2025), 173--200.
\href{https://doi.org/10.1007/s10589-025-00670-3}{doi:10.1007/s10589-025-00670-3}.

\bibitem{PWZ2006}
L. Pronzato, H. P. Wynn, and A. A. Zhigljavsky,
Asymptotic behaviour of a family of gradient algorithms in $\R^d$ and Hilbert spaces,
\emph{Mathematical Programming} \textbf{107} (2006), 409--438.
\href{https://doi.org/10.1007/s10107-005-0602-7}{doi:10.1007/s10107-005-0602-7}.
Author manuscript: \url{https://arxiv.org/abs/0802.4382v1}.

\bibitem{HDLZ2022}
Y. Huang, Y.-H. Dai, X.-W. Liu, and H. Zhang,
On the asymptotic convergence and acceleration of gradient methods,
\emph{Journal of Scientific Computing} \textbf{90} (2022), article 7.
\href{https://doi.org/10.1007/s10915-021-01685-8}{doi:10.1007/s10915-021-01685-8}.
Author manuscript: \url{https://arxiv.org/abs/1908.07111v1}.

\bibitem{He2025}
Y. He,
The worst-case root-convergence factor of GMRES(1),
2025 preprint.
\url{https://arxiv.org/abs/2501.10248v1}.

\bibitem{MMT2025b}
L. MacDonald, R. Murray, and R. Tappenden,
A Flexible Algorithmic Framework for Strictly Convex Quadratic Minimization,
2025 preprint, arXiv:2505.04047v1.
\url{https://arxiv.org/abs/2505.04047v1}.

\end{thebibliography}
\end{document}

